\scanpage{9}
\mspage{3}
Mais comment vérifier que $(p;g)$ définissent non seulement un endomorphisme, mais un automorphisme ? \unsure{Ça} semble délicat\ldots

Dans le cas de $\tau$, on a $p=-1$, et $g=l_1$ :
\[
\tau\rho=\operatorname{int}(l_1)\rho\tau
\]
On veut exprimer la commutation de $u$ (comme autom ext \inserted{\unsure{respectant} la structure à lacets}) non seulement à $\rho$ i.e. $\unsure{\mathfrak S_3^+|\mathcal S_3^+}$, mais à $\unsure{\mathfrak S_3|\mathcal S_3}$ tout entier. Il suffit d'exprimer la commutation mod autom extérieurs à $\sigma_\infty$, i.e. d'écrire qu'il existe $h\in\pi$ (également unique) tel que
\margin{\uline{NB} Ceci implique qu'il opère \unsure{triv\textsuperscript{t}} sur l'ens. des trois classes de conjugaison des ss-groupes de lacets --- ce qui précède s'applique.}
\[
u\sigma_\infty=\operatorname{int}(h)\deleted{\unsure{u}}\sigma_\infty u,
\]
et il suffira de l'exprimer pour les arguments $l_0$ et $l_1$. On trouve les conditions
\[
\underset{\substack{\Vert\\\underbrace{u(l_1)}_{\operatorname{int}(g)l_1^p}}}{u\sigma_\infty(l_0)}=\operatorname{int}(h)\underbrace{\sigma_\infty(l_0^p)}_{l_1^p}
\qquad\text{i.e.}\quad\operatorname{int}(g^{-1}h)l_1^p=l_1^p\quad\text{i.e.}
\]
\[
g^{-1}h\in L_1\quad\text{i.e.}\quad h=gl_1^{pq},\qquad q\in\hat{\mathbb Z}
\]
et
\[
\underbrace{u\sigma_\infty(l_1)}_{l_0}=\operatorname{int}(h)\sigma_\infty\underbrace{u(l_1)}_{\operatorname{int}(g)l_1^p}\qquad\text{ou}
\]
\[
l_0^p=\operatorname{int}(h\sigma_\infty(g))l_0^p\qquad\text{soit}
\]
\[
l_0^p=\operatorname{int}(gl_1^q\sigma_\infty(g))l_0^p
\]
ce qui équivaut à
\[
\deleted{\operatorname{int}(gl_1^q\sigma_\infty(g))L_1=L_0}
\]
i.e. la condition
\[
\exists_q\ \text{tel que}\quad gl_1^q\sigma_\infty(g)\in L_0
\]
\mspage{8}
